% Source: https://artofproblemsolving.com/community/c3h1194825_pascals_triangle_marathon
% Stable identity: aops-topic-1194825-q01
In Pascal's triangle, the $k$-th diagonal consists of the entries $\dbinom nk$ for $n=k,k+1,k+2,\dots$ (so diagonal $1$ is $1,2,3,\dots$ and diagonal $2$ is $1,3,6,10,\dots$). A student claims: ``If you compare two diagonals entry by entry in the same row and take the ratios, they form an arithmetic sequence.''
\begin{enumerate}
\item[(a)] Prove that the claim is true for two \emph{adjacent} diagonals: for fixed $k\ge0$, the ratios $\dfrac{\binom n{k+1}}{\binom nk}$, $n=k+1,k+2,\dots$, form an arithmetic sequence, and find its common difference.
\item[(b)] Show that the claim is false in general, by examining diagonals $1$ and $3$.
\end{enumerate}
